\chapter{Wstęp}\label{ch:wstep}

% ═══════════════════════════════════════════════════════════════════════════════
% Cel rozdziału: 3–5 stron.  Wprowadzić czytelnika w domenę, sformułować
% problem, zdefiniować cele i zakres, dać mapę rozdziałów.
% ═══════════════════════════════════════════════════════════════════════════════

\section{Motywacja}\label{sec:motywacja}

% TREŚĆ DO NAPISANIA — punkty do rozwinięcia w akapity:
%
% 1. Obciążenie administracyjne lekarzy w Polsce:
%    - Raport NIK / ZUS / Naczelna Izba Lekarska: lekarze POZ spędzają
%      ~40–50\% czasu pracy na dokumentacji, a nie na pacjencie.
%    - Wypalenie zawodowe lekarzy w Polsce (statystyki, ankiety).
%    - Problem szczególnie dotkliwy w POZ, gdzie wizyty trwają 10–20 minut,
%      a dokumentacja wymaga pełnych wpisów SOAP.
%
% 2. Wymogi regulacyjne w Polsce:
%    - Rozporządzenie MZ w sprawie rodzajów i zakresu dokumentacji medycznej
%      (Dz.U. 2020 poz. 666) — obowiązek prowadzenia dokumentacji.
%    - RODO (GDPR) — art. 9 (dane wrażliwe), art. 25 (privacy by design).
%    - Wymagania Narodowego Funduszu Zdrowia dotyczące kodowania procedur
%      i leków refundowanych.
%
% 3. Luka na rynku polskim:
%    - Istniejące rozwiązania (Nuance DAX, Abridge, Suki) — wyłącznie
%      anglojęzyczne, brak wsparcia dla polskiego systemu ochrony zdrowia.
%    - Brak zintegrowanego systemu CDSS dla języka polskiego z natywnym
%      wsparciem NFZ.
%    - Potencjał AI (STT, LLM, RAG) do odciążenia lekarzy.
%
% 4. Dlaczego teraz?
%    - Dojrzałość modeli STT dla języka polskiego (Whisper large-v3, WER <10\%).
%    - Dostępność lokalnych SLM (Phi, Llama, Bielik) z kwantyzacją na GPU
%      konsumenckie.
%    - Microsoft Presidio jako framework anonimizacji z możliwością
%      rozszerzenia o polskie recognizery (PESEL, NIP, imiona).

\section{Cel pracy}\label{sec:cel-pracy}

% TREŚĆ DO NAPISANIA — sformułowanie celu głównego + lista celów szczegółowych:
%
% CEL GŁÓWNY (1 zdanie):
%   Zaprojektowanie, implementacja i ewaluacja multimodalnego, wieloagentowego
%   systemu hybrydowego (CDSS), który automatyzuje tworzenie notatek SOAP
%   z nagrań wizyt lekarskich w języku polskim, z zachowaniem prywatności
%   danych (RODO) i integracją z bazą refundacji NFZ.
%
% CELE SZCZEGÓŁOWE (lista numerowana C1–C6):
%   C1. Opracowanie pipeline'u generacji syntetycznego datasetu dialogów
%       medycznych (scenariusze LLM + synteza mowy XTTSv2).
%   C2. Implementacja modułu transkrypcji (STT) z diaryzacją mówców
%       dla języka polskiego i ewaluacja jakości (WER/CER).
%   C3. Implementacja modułu strukturyzacji transkrypcji do formatu SOAP
%       przy użyciu lokalnego SLM.
%   C4. Implementacja modułu anonimizacji danych osobowych (Microsoft
%       Presidio + custom PL recognizery) z programistycznie wymuszaną
%       granicą zaufania (trust gate).
%   C5. Implementacja modułu chmurowego: analiza kliniczna LLM, RAG
%       i dopasowanie refundacji NFZ.
%   C6. Ewaluacja systemu end-to-end: jakość STT, poprawność SOAP,
%       skuteczność anonimizacji, trafność sugestii RAG.

\section{Zakres pracy}\label{sec:zakres}

% TREŚĆ DO NAPISANIA — co praca obejmuje vs. czego nie:
%
% OBEJMUJE:
%   - Projekt architektury wieloagentowej z podziałem on-premise/cloud.
%   - Implementacja prototypu (MVP) obejmującego pięć specjalizacji
%     medycznych: kardiologia, endokrynologia, neurologia, psychiatria,
%     ortopedia.
%   - Generacja syntetycznego datasetu (8 scenariuszy) jako ground truth.
%   - Ewaluacja na zbiorze syntetycznym z metrykami WER, CER, jakość SOAP
%     (ekspert kliniczny), precyzja/czułość anonimizacji, Recall@k RAG.
%   - Pulpit lekarza (frontend) jako interfejs akceptacji sugestii.
%
% NIE OBEJMUJE:
%   - Walidacja na prawdziwych danych klinicznych (wymaga zgody komisji
%     bioetycznej / IRB — poza zakresem czasowym pracy).
%   - Integracja z systemami HIS/EHR (np. mMedica, Asseco, Kamsoft).
%   - Certyfikacja jako wyrób medyczny (MDR).
%   - Pełna skalowalność produkcyjna (load balancing, HA, replikacja).
%   - Specjalizacje medyczne poza wymienionymi pięcioma.

\section{Struktura pracy}\label{sec:struktura}

% TREŚĆ DO NAPISANIA — przegląd rozdziałów, 1–2 zdania na rozdział:
Praca składa się z~siedmiu rozdziałów. W~rozdziale~\ref{ch:stan-wiedzy}
przedstawiono przegląd stanu wiedzy w~obszarach istotnych dla systemu
(rozpoznawanie mowy, modele językowe, anonimizacja PII). Rozdział~%
\ref{ch:wymagania} formalizuje wymagania funkcjonalne i~niefunkcjonalne
systemu. W~rozdziale~\ref{ch:architektura} opisano architekturę
wieloagentową z~podziałem na~strefę lokalną i~chmurową. Rozdział~%
\ref{ch:implementacja} omawia szczegóły implementacyjne poszczególnych
modułów. Rozdział~\ref{ch:ewaluacja} prezentuje wyniki ewaluacji systemu.
Rozdział~\ref{ch:wnioski} podsumowuje pracę i~wskazuje kierunki dalszego
rozwoju.

\section{Konwencje typograficzne}\label{sec:konwencje}

% Konwencje przejęte z szablonu PJATK:
W~tekście pracy stosowane są następujące konwencje:
\begin{itemize}
  \item \emph{kursywą} oznaczono terminy obcojęzyczne oraz nazwy własne;
  \item \code{czcionką o~stałej szerokości} oznaczono nazwy funkcji, pakietów,
        ścieżki plików i~fragmenty kodu;
  \item odwołania do~literatury zapisano w~nawiasach kwadratowych, np.~[1];
  \item rysunki, tabele i~listingi numerowane są niezależnie w~obrębie
        każdego rozdziału.
\end{itemize}
